\section{Requirements Analysis}

\subsection{Introduction}
This section outlines the requirements for the Smart Healthcare Management database system. The primary objective is to create a robust, scalable, and secure system that supports multiple user roles and complex medical data interactions.

\subsection{User Roles and Permissions}
The system will support the following user roles:
\begin{itemize}
    \item \textbf{Patient:} Can view their personal information, appointments, prescriptions, and lab results.
    \item \textbf{Doctor:} Can manage patient records, schedule appointments, issue prescriptions, and order lab tests.
    \item \textbf{Administrator:} Can manage hospital, department, and doctor information, overseeing the system's structural data.
\end{itemize}

\subsection{Core Entities}
The main entities identified for the system are:
\begin{itemize}
    \item Patient
    \item Doctor
    \item Hospital
    \item Department
    \item Appointment
    \item Prescription
    \item Lab Result
    \item Insurance
    \item Bill
\end{itemize}

\subsection{Core Relationships}
Key relationships between entities include:
\begin{itemize}
    \item A \textbf{Doctor} can have many \textbf{Patients} (1-to-N).
    \item A \textbf{Patient} can have many \textbf{Appointments} (1-to-N).
    \item A \textbf{Hospital} has many \textbf{Departments} (1-to-N), and each \textbf{Doctor} belongs to one \textbf{Department} (N-to-1).
    \item A \textbf{Doctor} can have multiple \textbf{Specialties} (N-to-M), implemented via a linking table.
    \item An \textbf{Appointment} must be associated with exactly one \textbf{Patient} and one \textbf{Doctor}.
\end{itemize}

\subsection{Constraints and Business Rules}
The system must enforce the following rules:
\begin{itemize}
    \item A prescription can only be written by an authorized doctor.
    \item Every appointment must be linked to an existing patient and doctor.
    \item Patient medical records are confidential and can only be accessed by authorized personnel.
\end{itemize}
